\addchap{Zusammenfassung}
\selectlanguage{ngerman}

Diese Arbeit untersucht die automatisierte Verifikation textueller
Benutzerschnittstellenanforderungen anhand geordneter visueller UI-Abläufe. Der
Ausgabevertrag kombiniert eine vierstufige Anforderungsentscheidung mit
expliziter Screenshot-Evidenz, Entscheidungen auf Claim-Ebene,
Unsicherheitsgründen und optionaler räumlicher Lokalisierung. Der Ansatz trennt
sichtbare Unterstützung von verborgenen Systemeigenschaften und erhält
Unsicherheit, wenn die verfügbaren visuellen Zustände keine verlässliche
positive oder negative Entscheidung erlauben.

Die empirische Untersuchung verwendet 258 geprüfte Verifikationsfälle aus 13
von Mind2Web abgeleiteten Web-Abläufen. Ein kontrolliertes Experiment variiert
automatische Claim-Zerlegung und Screenshot-Auswahl unabhängig voneinander.
Ergänzend werden Modellvergleiche, Wiederholungsläufe, ein
Flow-Cluster-Bootstrap, eine Robustheitsbedingung ohne verlässliche Reihenfolge
und eine qualitative Fehleranalyse durchgeführt. Im kontrollierten Vergleich
liefert die vollständige Screenshot-Sequenz die stärkste Grundlage für die
Labelvorhersage. Eine lexikalisch begrenzte Auswahl lässt entscheidende
Zustände aus und verschlechtert die Ergebnisqualität; die Wirkung der
Claim-Zerlegung hängt von der verfügbaren Evidenz ab. Explizite
Evidenzrepräsentationen verbessern weiterhin die Nachvollziehbarkeit, da sie
Fehler bei Auswahl, Interpretation und Aggregation getrennt sichtbar machen.

Die Evaluation verwendet einen geprüften Benchmark mit 13 Abläufen und
bezieht sich auf das, was in aufgezeichneten Screenshots sichtbar ist.
Verborgenes Backend-Verhalten wird nicht abgeleitet; die Ergebnisse
charakterisieren dieses Evidenzregime und nicht sämtliche UI-Entwürfe und
Implementierungen.

\selectlanguage{american}
